% Expected Outcomes and Innovations
\section{Expected Outcomes and Innovations}
\label{sec:outcomes}

This research is expected to deliver three artefacts and four scientific contributions, each
described in turn.

\subsection{Primary Deliverables}

The first deliverable is the proposed metric itself: a no-reference long-range video
consistency composite for super-resolution evaluation, with five reliability-gated
sub-metrics covering appearance, multi-scale temporal consistency, identity preservation,
colour-distribution stability, and colour-trajectory drift. The composite is decomposable:
the per-sub-metric scores and reliabilities are reported alongside the composite so that
each video's evaluation is traceable to the sub-metric driving it. The composite is fully
specified by a small set of canonical hyperparameters; alternative hyperparameter choices
can be explored through cached per-video sub-metric values without rescanning videos.

The second deliverable is a parameterised synthetic test set for long-video super-resolution
consistency assessment, covering four artefact families --- slow colour drift,
chunk-boundary jumps, periodic flicker, and identity degradation --- across multiple
severity levels per base video. Each generator is verified to produce zero output at zero
severity (identity within floating-point tolerance), monotonic per-frame statistics under
severity sweep, and seed-reproducible construction. The synthetic test set is the first
publicly characterised consistency benchmark targeting long-video super-resolution
specifically; existing video super-resolution benchmarks (REDS, Vimeo-90K, Vid4, UDM10,
SPMCS) target per-frame fidelity on short clips. The test set is extensible: additional
base videos, finer severity grids, and new artefact families can be added without
disturbing the existing infrastructure.

The third deliverable is the empirical-characterisation contribution of
Section~\ref{sec:problem}. The four findings --- the regime-shift case, the temporal-scale
crossover, the colour-drift blind spot of the baseline metric suite, and the
identity-collapse pathology of slow-fast pooling under heavy degradation --- are
documented, each with an identified root mechanism. To our knowledge these characterisations
have not previously been compiled into a single systematic study of long-video
super-resolution evaluation.

\subsection{Scientific Innovations}

The research introduces four innovations.

\textbf{Reliability-weighted softmax-log-mean composition as a design pattern for
no-reference long-video metrics.} The composition assigns weights through a sharp softmax
over per-sub-metric reliabilities, then computes a geometric mean of scores under those
weights. Sub-metrics whose reliability gates engage on a given video are downweighted
automatically, without per-sub-metric threshold-based hard exclusion. The composition is
sub-metric-agnostic: additional sub-metrics can be added as new failure modes are
characterised without disturbing the existing composition. To our knowledge this is the
first long-video super-resolution consistency metric of this form.

\textbf{Goodness-of-fit-gated reliability for trajectory-style sub-metrics.} Sub-metric E
addresses the colour-drift blind spot by fitting a linear regression to per-frame
Lab-channel means over the whole video, scoring an exponential decay in the maximum
absolute slope, and gating reliability by the maximum coefficient of determination across
the three channels. A drifting video produces a high $R^2$ (the slope is real, so the
metric speaks confidently); a flickering video produces $R^2$ near zero (sinusoidal
residuals dominate, so the metric correctly abstains). The pattern --- gate reliability by
goodness-of-fit, not by score magnitude --- is reusable for any slow-signal-versus-noisy-baseline
detection problem in video evaluation.

\textbf{Multi-scale temporal aggregation with configurable weighting.} The proposed
sub-metric T aggregates tOF over $k \in \{1, 5, 10, 30, 60, 120\}$; the per-video per-$k$
warping errors are cached so that alternative aggregation weightings can be applied at
composition time without recomputing the underlying optical flow. This enables principled
hyperparameter sweeps over the temporal-scale axis that produced the crossover finding of
Section~\ref{sec:problem:tof}.


